%%
%% trilogy_dag.tex
%% Statement-level logical DAG for Papers XVIII, XIX, XX.
%% Every node is a named proposition; every edge is an explicit
%% logical dependency (the conclusion requires the source).
%%
\documentclass[12pt,a4paper]{amsart}
\usepackage{amsmath,amssymb,amsthm,mathtools,enumitem,booktabs,hyperref}
\usepackage[T1]{fontenc}

\newtheorem{theorem}{Theorem}[section]
\newtheorem{lemma}[theorem]{Lemma}
\newtheorem{proposition}[theorem]{Proposition}
\newtheorem{corollary}[theorem]{Corollary}
\theoremstyle{definition}
\newtheorem{definition}[theorem]{Definition}
\newtheorem{hypothesis}[theorem]{Hypothesis}
\newtheorem{remark}[theorem]{Remark}

\newcommand{\R}{\mathbb{R}}
\newcommand{\normh}[2]{\|#1\|_{#2}}
\newcommand{\cA}{\mathcal{A}}
\newcommand{\cAg}{\mathcal{A}_\gamma}
\newcommand{\cK}{K_{B_c}}
\newcommand{\cKA}{K_{\cA}}

\begin{document}

\title{Logical DAG: Trilogy XVIII--XX (Statement Level)}
\date{May 2026}
\maketitle

\tableofcontents

%% ================================================================
\section{Purpose}
%% ================================================================

This document records the complete statement-level logical
dependency graph (DAG) for Papers~XVIII, XIX, XX.
Each node is a named result; each directed edge
$A\to B$ means ``$B$ logically requires $A$ in its proof''.
The DAG is organised by paper and by logical layer.

\medskip
\textbf{Notation.}
We write $[\text{XVIII}]$, $[\text{XIX}]$, $[\text{XX}]$
for the paper in which a result first appears.
Nodes marked (H) are hypotheses (not proved in the trilogy).
Nodes marked (O) are open.

%% ================================================================
\section{Layer 0: Foundations (agnostic to $\gamma$)}
\label{sec:layer0}
%% ================================================================

These results hold for any operator in $\mathfrak{T}_\gamma$
without the $\gamma$-sensitive hypotheses.

\medskip
\begin{center}
\renewcommand{\arraystretch}{1.35}
\begin{tabular}{l l l p{4.0cm}}
\toprule
\textbf{ID} & \textbf{Label} & \textbf{Paper}
  & \textbf{Dependencies} \\
\midrule
F1 & Algebraic split
  $S_K-\cKA = A_K - B_K$ & XVIII
  & Hilbert--Schmidt theory \\
F2 & $A_K = A_K^{\rm lin}+A_K^{\rm quad}$
  (tensor decomposition) & XIX
  & F1 \\
F3 & Tensor norm
  $\|f\otimes g\|\leq M\|f\|_{C^0}\|g\|_{C^0}$ & XVIII
  & Standard \\
F4 & Bessel tail bound
  (\texttt{lem:besseltail}) & XVIII
  & Bessel asymptotics \\
F5 & Tonelli/DCT (\texttt{rem:dct}) & XVIII
  & Standard \\
F6 & Slepian tail:
  $\|B_K\|\leq C_2^{1/2}e^{-\alpha K/2}$ & XVIII
  & Slepian prolate theory \\
F7 & Sequential $(c,K)$-limit & XVIII
  & F5, F6 \\
F8 & Kato norm-resolvent perturbation & XVIII
  & Kato VI \\
F9 & Total asymptotic order $\mathcal{S}$
  (Def.~\texttt{def:scale}) & XX
  & Standard (Hardy--Littlewood) \\
\bottomrule
\end{tabular}
\end{center}

%% ================================================================
\section{Layer 1: Airy-specific (XVIII--XIX)}
\label{sec:layer1}
%% ================================================================

These results are proved for the Airy case $V=\xi$,
$\gamma=2/3$, $\beta=1/3$.

\medskip
\begin{center}
\renewcommand{\arraystretch}{1.35}
\begin{tabular}{l l l p{4.5cm}}
\toprule
\textbf{ID} & \textbf{Label} & \textbf{Paper}
  & \textbf{Dependencies} \\
\midrule
A1 & Airy spectral growth:
  $a_j\sim(3\pi j/2)^{2/3}$ & XVIII
  & Classical Airy zeros \\
A2 & Langer approximation:
  $u_j^{(c)}\approx\mathrm{Ai}(c^{2/3}\zeta_j(\xi))$ & XVIII
  & Langer (1949), Olver (1954) \\
A3 & Langer $K$-dependence:
  $C_j=O(j^{2/3})$ (\texttt{rem:langerkdep}) & XVIII
  & A1, A2, Olver Ch.~12 \\
A4 & Turning-point coercivity window
  $c^{-1/3}$ & XVIII
  & A2, A3 \\
A5 & Qualitative BR3:
  $\cK\to\cKA$ in $L^2$ & XVIII
  & F1--F8, A2--A4 \\
A6 & Lin/quad split quantitative:
  $C_{\rm lin}=O(K^{5/3})$, $C_{\rm quad}=O(K^{7/3})$ & XIX
  & F2, F3, A3 + summation \\
A7 & $K_{\rm opt}=2/(5\alpha)\log c$ (Airy) & XIX
  & A6, saddle balance (Rem.~\ref{rem:kopt_airy}) \\
A8 & Airy rate:
  $O(c^{-1/3}(\log c)^{5/3})$ & XIX
  & A6, A7, F6, F9 \\
\bottomrule
\end{tabular}
\end{center}

%% ================================================================
\section{Layer 2: Universal (XX, conditional)}
\label{sec:layer2}
%% ================================================================

These results hold for $\cAg\in\mathfrak{T}_\gamma$
conditionally on Hypotheses H1--H2.

\medskip
\begin{center}
\renewcommand{\arraystretch}{1.35}
\begin{tabular}{l l l p{4.5cm}}
\toprule
\textbf{ID} & \textbf{Label} & \textbf{Paper}
  & \textbf{Dependencies} \\
\midrule
H1 & Pointwise Langer scaling law:
  $\|e_j\|\leq C_0 j^\gamma c^{-\beta}$
  & XX (H)
  & Verified: A3 (Airy);
    sketch: \texttt{olver\_growth\_lemma.tex};
    open: general \\
H2 & $C_\infty=\sup_j\|u_j\|_{C^0}<\infty$
  & XX (H)
  & Verified: Airy ($L^\infty$ bound);
    open: general \\
U1 & Summation:
  $\sum_{j=1}^K j^{n\gamma}\sim K^{1+n\gamma}/(1+n\gamma)$
  & XX
  & Euler--Maclaurin \\
U2 & Universal truncation geometry:
  $C_{\rm lin}=O(K^{1+\gamma})$,
  $C_{\rm quad}=O(K^{1+2\gamma})$
  & XX
  & F2, F3, H1, H2, U1 \\
U3 & Asymptotic order of layers
  (Prop.~\texttt{prop:asymp\_order})
  & XX
  & F9 \\
U4 & Layer-Survival Selection Principle
  (Thm.~\texttt{thm:layer\_survival})
  & XX
  & U2, U3 \\
U5 & Saddle balance:
  $K_{\rm opt}\sim 2\beta/(\alpha(1+\gamma))\log c$
  (first-order asymptotic)
  & XX
  & F6, U2 \\
U6 & Universal two-scale rate
  (Thm.~\texttt{thm:universal\_rate})
  & XX
  & F1, F6, U2 \\
U7 & Universal logarithmic rate:
  $O(c^{-\beta}(\log c)^{1+\gamma})$
  (Cor.~\texttt{cor:universal\_lograte})
  & XX
  & U4, U5, U6 \\
\bottomrule
\end{tabular}
\end{center}

%% ================================================================
\section{Layer 3: Open Problems}
\label{sec:layer3}
%% ================================================================

\begin{center}
\renewcommand{\arraystretch}{1.35}
\begin{tabular}{l l l p{4.0cm}}
\toprule
\textbf{ID} & \textbf{Label} & \textbf{Closes} & \textbf{Required tool} \\
\midrule
O1 & Prove H1 from Olver for general $V$
  & H1 $\to$ U2,U4,U6,U7
  & Olver-Watson expansion;
    regular variation of $V$ \\
O2 & Verify H2 for $V=|\xi|^\nu$
  & H2 $\to$ U2
  & WKB normalization bound \\
O3 & Lower bound $\Omega(K^{1+\gamma})$
  for $C_{\rm lin}$
  & sharpness of U4
  & spectral counting argument \\
O4 & Complete Olver-Watson expansion
  in \texttt{olver\_growth\_lemma.tex}
  & sketch $\to$ proof of H1
    (power family)
  & full Olver Ch.~12 residue \\
\bottomrule
\end{tabular}
\end{center}

%% ================================================================
\section{Critical Path to Full Theorem}
\label{sec:critical_path}
%% ================================================================

The minimal path from current state to
``U7 is an unconditional theorem for the power family''
is:
\[
  \text{O4} \Rightarrow \text{H1 (power family)}
  \Rightarrow \text{U2} \Rightarrow \text{U4}
  \Rightarrow \text{U7 (power family)}.
\]
O2 is independent of this path (it handles $C_\infty$,
which is expected to be standard but is currently unchecked
for $\nu\neq 1$).

For full generality beyond the power family,
O1 requires a \emph{regular variation} condition on $V$
(formally: $\xi V'(\xi)/V(\xi)\to p$ as $\xi\to\infty$
for some $p>0$), which would extend H1 to
all regularly varying potentials.

%% ================================================================
\section{Dependency Graph (Textual)}
\label{sec:textual_dag}
%% ================================================================

\begin{verbatim}
[Classical] Olver Ch.12, Langer 1949
   |                    |
   v                    v
  [A2] Langer approx   [A1] Airy zeros
   |                    |
   +---------+----------+
             |
             v
            [A3] C_j = O(j^{2/3})   [F2-F8]
             |                         |
             v                         |
            [A5] BR3 (XVIII)           |
                                       |
  [H1] j^gamma scaling   [H2] C_inf   |
   |                         |         |
   +----------+--------------+         |
              |                        |
              v                        |
             [U1] Euler-Maclaurin sum  |
              |                        |
              v                        |
             [U2] C_lin, C_quad        |
              |                        |
   [F9] Scale S           +------------+
     |         |           |
     v         v           v
    [U3] order [U5] K_opt [U6] two-scale bound
     |         |
     v         v
    [U4] Layer-Survival Theorem
     |         |         |
     +---------+---------+
               |
               v
             [U7] O(c^{-beta}(log c)^{1+gamma})  ■
\end{verbatim}

%% ================================================================
\section{Freeze Checklist}
\label{sec:freeze}
%% ================================================================

\begin{center}
\renewcommand{\arraystretch}{1.3}
\begin{tabular}{l l}
\toprule
\textbf{Item} & \textbf{Status} \\
\midrule
All F1--F9 nodes present in XVIII or XX & \checkmark \\
All A1--A8 nodes proved in XVIII--XIX & \checkmark \\
H1, H2 explicitly stated as hypotheses & \checkmark \\
U1--U7 proved conditional on H1, H2 & \checkmark \\
O1--O4 isolated and not blocking freeze & \checkmark \\
$K_{\rm opt}$ marked as first-order saddle & \checkmark \\
$c$-scale comparisons via Definition~\ref{def:scale} & \checkmark \\
``Conditional theorem'' status in abstract & \checkmark \\
\hline
O4 (Olver-Watson expansion) complete & open \\
O2 ($C_\infty$ for $\nu\neq 1$) verified & open \\
\bottomrule
\end{tabular}
\end{center}

\medskip
\textbf{Freeze verdict.}
XX is ready to freeze as a \emph{conditional theorem}.
The two open items (O2, O4) are self-contained;
they do not affect the internal consistency of XX
or the correctness of the proofs given H1--H2.

\end{document}
